\newcommand{\StudyVersion}{1.0 · modeled confirmation}

\newcommand{\ResultExecutive}{The model supports a strong conditional service business, but not a confirmed retained-Alpha flywheel. In the central case, the selected enterprise sells a useful module for an average \$25,599 against a scope-matched human-team reference of \$54,729, at \$4,672 fully allocated resource cost per useful outcome. Median two-year pre-tax profit is \$454,511 from \$50,000 seed capital. These are simulated outcomes under assumed task success, buyer value and demand.

After equalizing commercial policy choices, the selected enterprise uses competent Beta without retention in all three scenarios. All 4,608 fresh confirmation histories choose zero capability campaigns; no acquired Alpha, peer reuse or expanded offering is demonstrated there. Conservative assumptions produce losses in every world. A selected development trace runs the entire paid-work-to-peer-reuse lifecycle, but its complete economic return is negative. Profitable delivery and productive network growth remain separate findings.}

\newcommand{\ResultAbstract}{Can demanding paid work finance capabilities that create additional value beyond competent frontier intelligence? We construct a finite-cash, finite-capacity simulation of an autonomous optimization foundry supplying formally specified scheduling and planning policy modules. Human-team estimates use reported US service rates and explicitly assumed scope-matched effort; task performance, demand, buyer value and retained-capability effects remain exposed assumptions. Customer work A already pays for its formulation, learning and delivery. Retention B charges only additional generalization, assessment and lifecycle work; optional research C is separate.

A frozen policy screen, a disclosed pre-confirmation fairness amendment and 128 fresh worlds per scenario distinguish selected enterprise, same-price Beta, local Alpha and network Alpha. The central selected-enterprise mean pre-tax profit is \$497,058, with 127/128 profitable worlds; favorable assumptions give \$2,423,231 and 128/128, while conservative assumptions lose \$28,573 on average. The selected enterprise declines retention in all scenarios. No capability campaign occurs in any primary confirmation arm, so no causal sharing benefit or newly feasible Alpha offering is established. A separate screen saves \$3.04 per independent invention campaign but fails full lifecycle recovery under its capacity limit. The result is a conditional commercial surplus engine with explicit limits on the stronger network-growth claim, not evidence of actual sales, autonomous invention or realized improvement.}

\newcommand{\ResultProtocol}{The base protocol was frozen on 26 October 2026 before scientific policy search (SHA-256 850e6978a3946f8f9d01dd0c6920b49916e955d12b420a65c1241d3c5cd675f6). It uses 104 service weeks plus 24 runoff weeks and 24 development seeds per scenario. Its original catalogue contained 40 network-enabled policies and five fixed-price Beta policies. The score is mean discounted accrued profit plus one-quarter of the mean of the worst 10\% of development outcomes. The declared development insolvency limit is 5\%, with an explicit least-risk fallback if no policy qualifies.

After the original development screen, and before any confirmation seed was inspected, an identified comparison asymmetry was corrected in a registered amendment: Beta received the same price, scarcity-response and reserve choices as the network. The original protocol, training outputs and selection remain preserved. Twenty symmetric Beta policies were evaluated on the same development seeds, giving 60 unique final policies across the 40 network and 20 Beta choices. Network winners and all numerical model assumptions were unchanged. This amendment strengthens the comparator; it is disclosed because development results had already been seen.

The final policy selection was locked before 128 fresh confirmation seeds per scenario. Enterprise, separately selected Beta, network-enabled, same-price Beta, local Alpha and the prespecified ablations are evaluated on common world seeds. The selected enterprise chooses Beta-only in all three scenarios after the symmetry amendment. This is a genuine optional-investment result, not an omitted failed succession: the network-enabled alternative remains separately reported and is never relabeled as the optimizer's choice. Every searched policy enters the venture; prelaunch no-entry is outside the optimized catalogue. The no-separate-research ablation is not a separately optimized B-only controller.

One-factor stresses hold the selected central network policy fixed. The capacity study first compares 18 configurations on 12 separate development seeds, locks the chosen configuration at each demand level, and then uses 12 untouched capacity-confirmation seeds. Those holdout outcomes are not pooled with the capacity-selection data. A separate 256-campaign experiment tests retained screening and its exact zero-increment control.}

\newcommand{\ResultMainResults}{The positive commercial case is ordinary competent-Beta delivery. Under central assumptions, the selected enterprise completes 24.46 useful outcomes per world over 104 service weeks. Across the identical useful-delivery cohort, the human-team estimate averages \$54,729.08, the modeled bid \$25,599.20 and complete accrued resource expense \$4,672.21 per useful outcome. This allocation includes failed work, acquisition, common operations, setup and the historical allowance, including expense from zero-useful worlds. Mean pre-tax profit is \$497,058.08; the median is \$454,511.24 and the 5th--95th percentiles are \$154,703.34--\$964,240.36. One of 128 worlds loses money.

Favorable assumptions produce 47.66 useful outcomes per world. The matched human reference is \$88,702.85, bid \$55,896.30 and complete cost \$4,587.98. Mean pre-tax profit is \$2,423,230.89; median \$2,446,298.60 and 5th--95th percentiles \$1,562,475.99--\$3,137,891.09. All 128 worlds are profitable. These large profits follow the assumed valuable customer population and service-price reference; they are not measured demand or forecasts for an actual enterprise.

Conservative assumptions do not support launch. Only three useful outcomes occur across all 128 selected-enterprise worlds, and every world loses money. Mean loss is \$28,572.69, median loss \$28,244.12. Pooling the costs of all these attempted businesses over three useful deliveries yields \$1,234,391.29 per useful outcome. That extreme allocation reflects absent demand, not a million-dollar generation call. The searched policies had already entered the venture; a prelaunch no-entry alternative was outside the policy catalogue.

No primary world becomes insolvent under the \$50,000 seed and stop/reserve rules. Zero events in 128 worlds has a Wilson 95\% upper bound of 2.91\% under this model; it is neither certainty nor protection against parameter misspecification. At the central and favorable medians, available cash after obligations and reserves is \$391,973.65 and \$1,925,649.92 respectively. The separately recorded earned portions are \$341,973.65 and \$1,875,649.92; available balances are not all free operating profit.}

\newcommand{\ResultPricing}{The declared base-price ratios are 0.35, 0.50, 0.65, 0.80 and 0.95 of the task-matched human reference, combined with permitted scarcity and reserve responses. The final selected Beta policies use base ratios 0.50, 0.50 and 0.65 in conservative, central and favorable scenarios. Scarcity responses and the mix of useful modules explain why actual cohort-average bids differ from base-reference multiplication.

Among contacted prospects, the selected enterprise obtains a modeled AI choice in 0.37\%, 25.45\% and 41.24\% respectively, before support and capacity filtering. The human option, postponement and decline remain available. These are scenario implications of the utility assumptions, not calibrated elasticities. Substantial discounts were economically preferred within this catalogue; the study does not establish that a five- or ten-percent discount would attract the same buyers.

A service price floor and the buyer ceiling must overlap. The appendix reports fixed-mix break-even arithmetic separately from executed price policies. Changing price can also change trust-sensitive selection, completion beliefs, queueing and failure mix, so those algebraic thresholds cannot be read as a fresh demand rerun.}

\newcommand{\ResultInvestment}{Paid work leaves modeled candidates even when retention is declined. The central network arm produces 7.88 candidates and 6.64 rights-cleared candidates per world on average; favorable values are 24.23 and 23.13. Nevertheless, no B retention campaign, optional C campaign, Alpha admission, transfer or revalidation occurs in any of the 4,608 primary confirmation histories. No already-charged source-job development cost is recharged.

The network controller still pays for automated opportunity assessment: mean B expense is \$13.31, \$362.50 and \$398.56 per world across the scenarios. These are consideration costs, not successful extraction or paid Alpha-admission trials. In the saved central and favorable first-world examples, paid proposals are affordable and bookable but stop at the controller because complete prospective evidence/upkeep and foregone delivery contribution dominate the expected benefit. A cheap extraction step alone is insufficient.

This is a result about the registered controller, resource reservations, finite candidate catalogue and forecasting assumptions. It does not prove that retaining knowledge is intrinsically uneconomic, nor does it require an asset to repay paid-job A twice. A separately optimized B-only controller was not searched. The predeclared no-separate-research diagnostic instead holds the chosen network controller fixed and is numerically identical to it in every primary confirmation world.}

\newcommand{\ResultNewOpportunities}{The primary confirmation establishes zero newly serviceable useful outcomes attributable to acquired Alpha. With no admitted asset, graph availability and recipient transfer cannot create a new offering. Ordinary paid candidates and higher ordinary throughput are not substituted for that missing endpoint.

The capacity study includes a rare selected development trajectory in which an admitted bundle makes seven useful jobs newly serviceable under the AI enterprise's budget and deadline. The same trajectory produces fewer total useful outcomes than its competent-Beta pair and loses money incrementally. It therefore illustrates a changed feasible set without establishing a net productive expansion.}

\newcommand{\ResultNetwork}{Network minus local profit and useful-outcome differences are exactly zero in every primary world: no new asset reaches either portfolio. This is absence of an activated sharing mechanism in the tested regime, not evidence that mandatory availability would be valueless after a useful admission. No-paid-learning differences are also zero because candidate emergence never leads to a committed retention campaign.

The network-minus-same-price-Beta mean pre-tax profit differences are -\$10.91 (Monte Carlo 95\% interval -\$21.95 to \$0.12), -\$2,309.36 (-\$5,157.02 to \$538.30), and \$1,739.81 (-\$17,519.10 to \$20,998.73). The favorable positive point estimate is not acquired-Alpha value: the interval spans zero and all portfolios contain zero new assets. The contrast includes automated consideration expense and the resulting capacity, pricing and delivery trajectory changes. The symmetry-amended enterprise choice remains Beta in every scenario.}

\newcommand{\ResultMovingBeta}{The assumed baseline upgrade changes ordinary resource requirements and can therefore improve service economics without retention. Network-minus-static-Beta mean profit is \$12,091.90 in the central case (Monte Carlo 95\% \$4,442.38--\$19,741.43) and \$64,189.74 in the favorable case (\$49,244.36--\$79,135.12). Mean useful-outcome gains are 0.63 and 1.37. These are consequences of the hypothetical future-Beta assumptions, not observed model releases or learned self-improvement.

Because the primary retained portfolios are empty, these runs cannot estimate realized erosion of an admitted asset. The code's version, freshness and retirement controls are independently exercised by the software qualification, and the selected development example retires when timely revalidation is unavailable. Static-Beta contrasts must not be labeled pure Alpha-decay estimates. The development witness retires in week 53, before its hypothetical Beta upgrade in week 78, so it does not establish economic survival or erosion of acquired Alpha through that upgrade.}

\newcommand{\ResultAdmission}{The local graph qualification receipt records 41 passing software controls. They cover direct-label identity fixtures, separated authorities, admission prerequisites, mandatory shareable discovery, freshness, moving-Beta revalidation, revocation, composition, encryption boundaries, failure memory and protected-material filtering. These are observed code checks. The economic probe outcomes and identities used by the experiment remain simulated, and semantic secrecy or live ENS/on-chain enforcement is not established.

The main empirical unit for economic comparisons is an independent network/world seed. A useful peer-use job is not itself causal evidence of sharing value; network minus local supplies that comparison. Prior-surplus-funded spending does not establish asset payback. The no-Alpha-effect arm removes the declared mechanisms globally, including initially endowed equivalents, and is a mechanism-null sensitivity rather than an acquired-asset-only ablation.}

\newcommand{\ResultWitness}{A selected capacity-development example, L144C8Q4 at seed 1035001, shows the implemented path without being fresh confirmation. Atlas delivers rolling-policy job j0 at a modeled invoice of \$96,325.82, collected in week 5. Its direct A bill is \$317.37, including workflow learning and support; acquisition, common operations and the historical allowance remain separate. A reusable representation/certificate candidate leads to week-8 B spending of \$17.91 extraction plus \$4,534.41 actual assessment, funded from prior surplus with zero A recharge. A conservative 31-week reservation precedes week-40 admission on 12 fresh paired clusters and a modeled lifetime lower net bound of \$5,652.90.

The bundle becomes available to all three eligible members. Orion lacked the equivalent, adapts it in week 44 for \$235.08 and completes four useful peer jobs among eight peer attempts; total uses are 23. The asset is technically productive in the model's latent state, which does not establish economic Alpha. Its positive admission forecast is an economic false positive ex post. It retires in week 53 when fresh revalidation is unavailable. Paired workflow resource savings of \$1,992.36 minus referenced B of \$5,480.41 and access of \$46 leave -\$3,534.05 before shared discovery. Whole-policy network minus same-price Beta is -\$281,949.96 and 15 fewer useful outcomes. Admission and useful peer work did not establish positive complete return.

At admission, the \$22,462.84 prospective remaining lifecycle required at least 420 uses at the \$53.53 per-use lower value, versus 525.2 forecast and 23 realized uses. A conditional 23-week minimum useful life assumes the same demand density and feasible revalidation; initial evidence lasted 13 weeks and that continuation condition failed. This one post-selection development illustration is never pooled with the 36 fresh capacity condition/world evaluations (12 shared fresh seeds), which had zero admissions. The 31-week serial blind reservation and forecast reuse without schedulable continuing revalidation are limitations of this assumed workflow/controller, not measured limits of frontier intelligence.}

\newcommand{\ResultCampaigns}{In 256 independent modeled campaigns with 96 candidate opportunities each, the family-specific screen reduced mean variable campaign expense from \$332.70 to \$329.66. The paired saving was \$3.04 per campaign (Monte Carlo 95\% interval \$2.84--\$3.24), approximately 0.914\%. This is a limited modeled improvement in experiment expense. The added screen targets failed mechanical evaluations after generation; it does not reduce the expensive candidate-generation token budget. This explains why a useful local efficiency gain has a small financial effect under these registered resource assumptions.

Useful validated outcomes were 131/256 under expert Beta and 132/256 with the screen. The paired useful-rate difference interval was -0.375 to 1.156 percentage points; no positive quality effect is established. Each arm had 21 local evaluation false passes: 16 invalid actual proposals plus five non-inventive fallback outputs. These are campaign-check outcomes, not graph Alpha admissions. A post-freeze audit found six fallback outputs unnecessarily reviewed per arm, costing \$81.35 per arm. The charges match case by case, so all paired savings, useful-outcome counts and the resource break-even are unchanged. The frozen results and this endpoint limitation are both preserved in the evidence companion. The exact zero-increment control matched the baseline on every registered comparison.

The complete separately built protocol lifecycle costs \$3,209.91, including full optional research and subsequent evaluation/upkeep. Resource-saving payback requires 1,056 independent future campaigns at the observed mean saving. The stated 65-week input quota supports at most 44 such campaigns before reserving capacity for ordinary paid work, yielding only \$133.80 gross expected savings. The full-cost acceleration claim therefore fails under this mechanism experiment. It does not justify forcing optional C onto a cheaper candidate whose development was already paid in A.

The screen's availability and additional information are specified hypotheses. This experiment executes candidate-outcome and cost simulation; it does not discover an arbitrary new algorithm, demonstrate a realized self-improving enterprise, or pay network distribution and admission costs by assumption.}

\newcommand{\ResultAutomation}{Every primary production, intake, evaluation, correction and business-control role is modeled as automated, with zero routine provider human hours. This is the supported-scope hypothesis, not a live operational achievement. Central and favorable selected-enterprise useful throughput averages 1.02 and 1.99 outcomes per month. The median of each world's median reserved delivery latency is two and three weeks; the median world's 95th-percentile queue wait is two weeks in both cases. Mean capacity refusals over the service horizon are 22.18 and 133.21. Unused worst-case retry reservations are retained until promised completion.

False passes remain despite automation. The average world-level false-initial-acceptance share among apparently accepted outputs is 2.77\% centrally and 1.40\% favorably; it is an assumed error-process result, not measured safety. Later-refunded cases are included; the summaries do not identify a residual unrefunded-error count. Scope filtering, correction, recovery and refunds are charged. Reserved quota utilization is distinct from actual billed API usage.

The independent capacity holdouts select eight slots and fourfold quotas at 12, 48 and 144 hypothetical prospects/month. Mean useful outcomes over 24 months are 64.50, 192.42 and 261.25, and mean full pre-tax profits are \$3.722M, \$11.399M and \$14.865M under favorable market assumptions. All 36 fresh capacity condition/world evaluations, using 12 shared fresh seeds across the three demand levels, have zero admissions and zero useful peer reuse. At the highest demand, network minus same-price Beta loses \$860,684 on average (Monte Carlo 95\% -\$1.276M to -\$445,545). More demand and compute did not establish a profitable retained network. Each demand condition has only 12 holdouts; zero insolvencies permits a 24.25\% Wilson upper bound.}

\newcommand{\ResultFragility}{The completed sensitivity sweep contains 136 registered settings, each paired across the same 20 fresh exogenous seeds: 2,720 parameter/world evaluations and 5,440 policy histories, not 2,720 independent replications. All records reconcile, including 39 repeated-configuration consistency controls. The central commercial policy is held fixed; these executed points do not reoptimize pricing or recalibrate demand. Zero events in 20 worlds still permits a Wilson 95\% upper rate of 16.1\%.

Commercial value and demand determine the strongest breakpoints. At the base human-team reference of \$7,167, mean network profit is -\$35,658 and no world is profitable. At \$14,334, mean profit is +\$13,931 but the median is -\$24,096, only 8/20 worlds are profitable and 2/20 become insolvent. At \$28,668 and \$57,337, means are +\$487,399 and +\$1,221,536, with 19/20 and 20/20 profitable worlds. Thus the positive mean changes sign between the two lowest executed reference points, but this is a discrete model bracket, not a measured market threshold. With zero prospects, the mean loss is \$40,908. Two prospects per month yield +\$104,247 on average but only 13/20 profitable worlds; 24 per month yield +\$662,224 and 20/20 profitable worlds. Buyer value is independent of the human-service reference and neither demand nor willingness to pay has been observed.

Cash and scarce resources can overturn a positive average. Initial seed of \$20,000 and setup cost of \$40,000 each cause insolvency in all 20 worlds. An input quota of three million tokens per week produces five insolvencies and mean profit of -\$21,072. A \$1,000 acquisition cost per contacted prospect produces three insolvencies despite mean profit of +\$202,717. Tenfold or fiftyfold token or compute usage leaves no useful delivered outcomes and mean losses of roughly \$35,600--\$36,500. These are resource-demand shocks that change both cost and quota/deadline feasibility, not isolated tariff-price tests and not evidence that frontier models cannot perform the work. Payment-delay results are nonmonotone because timing changes collection, capacity and buyer-feedback trajectories; no monotonic causal claim is imposed.

Across the sweep, nine capability campaigns appear in eight seed labels: four B retention attempts and five separate C research attempts. All nine fail final admission. The aggregate campaign fields include two technically productive candidates with planned admission, but no actual admitted asset; the summaries do not distinguish an economic-gate failure from stale-Beta or graph rejection at completion. The only campaign-triggering settings are 48 prospects per month and threefold service latency. No setting yields admitted Alpha, useful peer reuse or a positive lower Monte Carlo confidence bound on the network-minus-Beta mean. Ninety-six settings have a positive point estimate, but these arise without acquired assets and cannot establish network Alpha. Null candidate productivity, candidate-emergence and retention-cost changes often leave outcomes identical because the retention mechanism never activates. Such inactivity is not robustness of a successful mechanism.

Figure 5 and the complete source tables preserve adverse and null conditions. The separately locked capacity holdouts likewise contain no new admissions: increasing demand and hardware produces more ordinary useful work, while the highest-demand network policy loses \$860,684 relative to same-price Beta on average. The development witness and narrow campaign-efficiency gain therefore remain conditional illustrations with negative complete economic recovery, rather than successful cases hidden by an aggregate statistic.}

\newcommand{\ResultConclusion}{The strongest supported positive finding is a conditional autonomous service business: formal, valuable, machine-checkable module work can cover complete resource and business cost while selling below a matched multidisciplinary human-service estimate. The central and favorable scenarios retain substantial cash after obligations. This conclusion depends on the assumed buyer value, trust, recurrence, structured scope and workflow quality; external tariffs and published research do not verify those commercial assumptions.

The stronger retained-Alpha claim is not confirmed. Equal commercial policy freedoms select Beta-only in every scenario; primary confirmation commits no capability campaign. Higher-capacity holdouts do not recover the missing admission and peer-value result. The rare selected development admission demonstrates a functioning lifecycle and useful reuse, but loses money after its incremental costs and foregone work. The separate invention screen improves a narrow expense endpoint and fails full-cost recovery.

The study therefore identifies a boundary to test: paid work can fund inexpensive extraction, but evidence, adaptation, revalidation and scarce-capacity opportunity cost must fit within the candidate's remaining useful life and realizable additional value. Holding cash is the correct registered response when that inequality fails. These negative outcomes are retained rather than converted into frontier-model limitations or an assumed success loop. Later empirical work can test the stated task and market assumptions, and later separately registered research can investigate more efficient evidence protocols or a broader policy catalogue.}

\newcommand{\ResultFrench}{Le résultat positif est une entreprise de services autonome conditionnelle, pas une boucle Alpha rentable confirmée. Dans le scénario central, l'entreprise sélectionnée livre en moyenne 24,46 résultats utiles en deux ans. Pour le même ensemble de résultats, la référence humaine multidisciplinaire est de 54 729,08 USD, le prix proposé de 25 599,20 USD et le coût complet alloué de 4 672,21 USD. Le bénéfice médian avant impôt atteint 454 511,24 USD ; 127 mondes sur 128 sont rentables. Le scénario favorable donne un bénéfice médian de 2 446 298,60 USD. Ces chiffres proviennent de la simulation et de ses hypothèses de valeur, de demande, de qualité et de périmètre.

Les taux horaires humains sont publiés, mais les heures et les conditions de paiement à l'acceptation sont supposées. La valeur client est distincte du coût humain. Les clients peuvent choisir le service humain, différer ou refuser. Le scénario conservateur ne justifie pas le lancement : les 128 mondes perdent de l'argent, avec une perte médiane de 28 244,12 USD. L'absence d'insolvabilité observée sous un capital initial de 50 000 USD ne garantit ni la rentabilité ni la sécurité réelle.

Après une correction transparente du catalogue avant confirmation, Bêta seul est sélectionné dans les trois scénarios. Les 4 608 trajectoires de confirmation ne contiennent aucune campagne de conservation ou de recherche. Les candidats issus du travail rémunéré ne deviennent donc ni Alpha admis, ni avantage chez un autre agent, ni offre nouvellement réalisable. Le travail A paie déjà formulation, apprentissage et livraison. B ne comprend que le travail supplémentaire de conservation et de réutilisation ; C reste une recherche distincte facultative. A n'est jamais facturé deux fois.

Un exemple sélectionné dans les données de développement réalise le parcours complet : travail payé, conservation financée par surplus, admission, quatre travaux utiles chez Orion, puis retrait faute de revalidation disponible. Son rendement complet reste négatif : la politique réseau perd 281 949,96 USD face à Bêta au même prix. Il ne remplace pas les 36 mondes inédits de capacité, qui ont zéro admission. Un protocole expérimental distinct économise 3,04 USD par campagne, mais nécessite 1 056 campagnes pour couvrir son cycle complet, contre une capacité de durée de 44 avant les travaux ordinaires.

L'étude observe l'exécution du simulateur et des contrôles logiciels, pas des achats, des inventions autonomes ou des transferts réels. Elle préserve une conclusion précise : une prestation utile peut produire un surplus important ; la conservation Alpha exige encore que preuve, adaptation, entretien et travail déplacé soient couverts pendant une durée utile réalisable. Conserver les fonds ou abandonner un candidat est un résultat économique valable.}

\newcommand{\ResultAudit}{The completed primary confirmation audit covers 384 scenario/world evaluations and 4,608 policy histories, using 128 independent seeds per scenario. The same seed labels couple scenarios, so the 384 evaluations are not pooled as jointly independent observations. Its bookkeeping checks pass over 2,633,597 synthetic bills, 2,743,106 cash events and 183,833 modeled job records. Every primary arm has zero capability campaigns and zero revalidations. These are observed simulator/audit counts, not empirical production throughput. Actual customer jobs, collections, deployed agent authentications and frontier invention runs are zero.

The completed sensitivity audit covers 136 registered settings and 5,440 histories on 20 shared exogenous seed labels; all nine capability campaigns fail final admission. Capacity development and fresh holdout total 504 audited histories, kept separate by design. The separate invention experiment contains 256 independent paired campaigns. The source-level statistical-unit qualification documents which repeated seed labels couple scenarios and conditions, so history and event counts are not presented as independent commercial trials.

The local graph receipt records 41 passing controls, and relocation of the source packet passes all seven software-test programs. Source, protocol and final-policy hashes bind the analysis. A/B/C partitions reconcile exactly before display rounding, protected obligations precede commitment, source A is never recharged, and refunds are netted from receipts without duplicating resource expense. Monte Carlo intervals describe simulation sampling conditional on the model, not uncertainty in assumed real-world parameters.

A storage interruption paused and resumed the existing simulation processes without source, seed or numerical changes. Completed large ledgers are retained through lossless durable archives with byte/hash/count receipts. A zero-byte local placeholder is not the ledger; the restoration manifest identifies the complete archive. The source-only repository publication excludes binary artifacts, so exact full reproduction requires the companion packet, including the preserved frozen historical checkpoint.}

\newcommand{\ResultEconomicsFigure}{\studyfigure{../paper/figures/economics.pdf}{1}{Matched useful-outcome prices and full resource cost, alongside independent-world profit distributions. Full cost includes failed and zero-useful worlds; profit is not a subtraction of unpaired medians.}}

\newcommand{\ResultProfitFigure}{}

\newcommand{\ResultPricingFigure}{}

\newcommand{\ResultInvestmentFigure}{\studyfigure{../paper/figures/investment.pdf}{1}{Mean A/B/C/common expense and candidate-to-retention counts. Paid development is charged once. Apparent admission remains distinct from latent productive value.}}

\newcommand{\ResultCashFigure}{}

\newcommand{\ResultNetworkFigure}{\studyfigure{../paper/figures/network.pdf}{1}{Paired total economic effects on fresh confirmation worlds. Same-price Beta tests the retention-enabled policy; local tests network availability. No asset was acquired in primary confirmation. Intervals are conditional Monte Carlo mean intervals.}}

\newcommand{\ResultMovingBetaFigure}{}

\newcommand{\ResultCampaignFigure}{\studyfigure{../paper/figures/campaign.pdf}{1}{The separate experiment protocol reduces variable evaluation expense but fails full lifecycle recovery under the registered lifetime and quota assumptions. Campaign volume is not customer demand.}}

\newcommand{\ResultSensitivityFigure}{\studyfigure{../paper/figures/sensitivity.pdf}{1}{Prespecified fragility cases and fresh capacity holdouts. The 12-world capacity intervals remain broad and do not establish a low-risk scale guarantee.}}
